% 実測用の「本物の論文」。pdflatex で組む。
%
% 合成 PDF（make-paper-pdf.ts）は自分で座標を置くので、自分が想定した壊れ方しか
% 再現できない。実際に起きた崩れ（1 ページ目で左右の段が語単位に混ざる、付録が
% 丸ごと参考文献になる）は、どれも合成の 14 シナリオ全部 100% のまま素通りしていた。
% だからここは pdflatex に組ませる。
%
% 正解（ground truth）は **この .tex から機械的に** 取り出す（scripts/real-paper-truth.ts）。
% 規則はひとつ:「本文の 1 文は \s{...} で 1 行に書く」。
% PDF から取り出した結果を正解にしてはいけない — それでは何も測れない。
%
% \s は本文にそのまま展開されるだけなので、組版には一切影響しない。
%
% 書式の制約（正解抽出を素朴なままに保つため。詳細は real-paper-truth.ts）:
%   - \s{...} の中では、\cite / 単一文字の $x$ / ~ / \emph 系だけを使う。
%   - 図表番号は \ref を使わず literal に書く（\ref を解決する仕掛けを足さない）。
\documentclass[10pt,twocolumn,a4paper]{article}

% 本文を Times にする（IEEE / ACL の体裁）。標準 14 書体なので PDF にフォントが
% 埋め込まれず、リポジトリに入れる PDF が半分の大きさで済むという実利もある。
\usepackage{mathptmx}
\usepackage[margin=2cm,columnsep=10pt]{geometry}
\usepackage{amsmath}
\usepackage{fancyhdr}
\usepackage{booktabs}

% 本文 1 文を包むだけのマクロ（正解抽出の目印）。
\newcommand{\s}[1]{#1}

% 柱とノンブル。arXiv のプレプリントでよく見る形。
\pagestyle{fancy}
\fancyhf{}
\fancyhead[L]{Preprint. Under review.}
\fancyhead[R]{Block Sparse Attention}
\fancyfoot[C]{\thepage}
\renewcommand{\headrulewidth}{0.4pt}

\title{Block Sparse Attention for Long Documents}
\author{
  Jane Q. Doe\thanks{Work done during an internship at Seaside Corporation.}\\ Department of Computer Science, Northern Institute of Technology\\ \texttt{jane.doe@example.edu}
  \and
  Alan R. Roe\\ Center for Language Systems, Riverside University\\ \texttt{alan.roe@example.edu}
  \and
  Mei Tanaka\\ Advanced Research Laboratory, Seaside Corporation\\ \texttt{mei.tanaka@example.com}
}
\date{}

\begin{document}

% 段抜きの表題とアブストラクト（NeurIPS / CVPR と同じ組み方）。
\makeatletter
\twocolumn[
  \begin{@twocolumnfalse}
  \maketitle
  \begin{abstract}
  \s{Transformer models spend most of their compute on attention when the input is long.}
  \s{We introduce a block sparse attention kernel that keeps the accuracy of dense attention at a fraction of its cost.}
  \s{The kernel selects a small set of key blocks for every query block and evaluates attention only inside that set.}
  \s{Across five long document benchmarks our method matches the dense baseline within one point of accuracy.}
  \s{It runs three times faster at a sequence length of thirty two thousand tokens.}
  \s{Longer contexts also change how practitioners build systems, because the retrieval step no longer has to be precise.}
  \s{We show that the selection pass can be shared across layers with almost no loss in accuracy.}
  \s{We further show that the pattern learned by the selection pass is stable across random seeds.}
  \s{A qualitative analysis suggests that the retained blocks correspond to the passages a human annotator would mark as relevant.}
  \s{The kernel is implemented once and reused for both training and inference.}
  \s{Training with the sparse kernel from scratch reaches the same loss as dense training after the same number of steps.}
  \s{We release the kernel and the trained checkpoints so that the results can be reproduced.}
  \end{abstract}
  \vspace{1em}
  \end{@twocolumnfalse}
]
\makeatother

\section{Introduction}

\s{Retrieval augmented systems concatenate dozens of documents into a single prompt, and agentic systems accumulate context over many turns.}
\s{Both patterns push sequence lengths far beyond what dense attention was designed for.}
\s{Attention computes a score for every pair of positions, so its cost grows with the square of the sequence length.}
\s{At thirty two thousand tokens the attention matrix alone occupies more memory than the model weights.}
\s{Practitioners respond by truncating the context, which silently discards the evidence the model was supposed to read.}

\s{A long line of work replaces the dense matrix with a sparse one.}
\s{Sliding window attention keeps a fixed band around the diagonal and drops everything else.}
\s{Random attention samples a fixed number of keys for every query.}
\s{Both schemes are cheap, but they fix the sparsity pattern before seeing the data.}
\s{When the evidence a query needs sits outside the pattern, the model cannot recover it.}

\s{We take a different route and let the model choose the pattern at run time.}
\s{Queries and keys are grouped into contiguous blocks of sixty four positions.}
\s{A cheap scoring pass compares block summaries and keeps the highest scoring key blocks for each query block.}
\s{The expensive pass then runs ordinary attention inside the selected blocks only.}
\s{Because the selection is made from the data, the pattern adapts to each input.}

\s{Our contributions are the following.}
\s{We describe a block selection rule that costs a small fraction of a full attention pass.}
\s{We prove that the approximation error is bounded by the mass of the discarded blocks.}
\s{We give a fused kernel that reaches eighty percent of the peak throughput of dense attention on the same hardware.}
\s{We evaluate the method on five benchmarks and report both accuracy and wall clock time.}

\section{Background}

\s{We write the query, key and value matrices as Q, K and V.}
\s{Each has one row per position and one column per head dimension.}
\s{Dense attention forms the full score matrix and normalises it with a softmax over the key axis.}
\s{The result is a convex combination of value rows, one combination per query.}

\s{The scaled dot product attention that we start from is written as Equation 1.}
\begin{equation}
  \mathrm{Attention}(Q, K, V) = \mathrm{softmax}\!\left( \frac{Q K^{\top}}{\sqrt{d}} \right) V ,
  \label{eq:attn}
\end{equation}
\s{Here $d$ is the head dimension and the softmax runs along the key axis.}

\s{When the sequence length grows, the softmax denominator grows with it, and the sum is dominated by a large number of near zero terms.}
\s{This observation is the reason sparse attention works at all.}
\s{If most of the probability mass sits on a few keys, the rest can be dropped without changing the output much.}
\s{The difficulty is finding those keys without computing the scores first.}

\section{Method}

\s{We partition the query positions into blocks of equal size and do the same for the key positions.}
\s{For each block we keep a summary vector, defined as the mean of the rows in that block.}
\s{The selection pass computes the dot product between query summaries and key summaries.}
\s{This is a small matrix, smaller than the full score matrix by the square of the block size.}
\s{For every query block we keep the top scoring key blocks and discard the rest.}

\s{The selection score between query block $i$ and key block $j$ is given by Equation 2.}
\begin{equation}
  s_{ij} = \frac{\bar{q}_i^{\top} \bar{k}_j}{\sqrt{d}} ,
  \label{eq:sel}
\end{equation}
\s{The retained set is then the arg max of that score over the second index.}

\s{The retained set always includes the diagonal block and the first block of the sequence.}
\s{The diagonal block carries local context, which every position needs.}
\s{The first block carries the instruction prefix, which queries attend to throughout the sequence.}
\s{Adding these two blocks by hand costs nothing and removes the most common failure of learned patterns.}

\subsection{Error Bound}

\s{The error introduced by dropping a block is bounded by the probability mass that the block would have received.}
\s{Because the softmax is one Lipschitz in the logits, the output error is at most twice that mass.}
\s{In practice the retained blocks hold more than ninety nine percent of the mass at every layer we inspected.}

\subsection{Kernel Design}

\s{A naive implementation gathers the selected blocks into a dense buffer and calls a standard kernel.}
\s{The gather costs as much as the attention it replaces, so we fuse the two steps.}
\s{Our kernel walks the selected blocks directly and keeps the running softmax statistics in registers.}
\s{Memory traffic is therefore proportional to the number of retained blocks rather than to the sequence length.}

\begin{table}[t]
  \centering
  \begin{tabular}{lrr}
    \toprule
    Model & Accuracy & F1 \\
    \midrule
    Dense attention & 81.2 & 79.4 \\
    Sliding window & 76.5 & 74.1 \\
    Block sparse (ours) & 88.7 & 86.1 \\
    \bottomrule
  \end{tabular}
  \caption{Accuracy on the held out split of the long document benchmark. Higher is better in both columns.}
  \label{tab:main}
\end{table}

\section{Experiments}

\s{We evaluate on five long document benchmarks that range from eight thousand to sixty four thousand tokens.}
\s{All models share the same backbone and differ only in the attention implementation.}
\s{We report accuracy on the held out split and wall clock time for a single forward pass.}
\s{Every number is the mean of three runs with different seeds.}
\s{We then report the main results in Table 1.}

\s{The dense baseline is the strongest of the fixed pattern methods but also the slowest.}
\s{Sliding window attention is fast and clearly worse, which matches earlier reports \cite{doe2024sparse}.}
\s{Our method is both faster than the dense baseline and more accurate than it.}
\s{The gain in accuracy is small but consistent across the five benchmarks.}

\begin{figure}[t]
  \centering
  \fbox{\rule{0pt}{3.2cm}\rule{6cm}{0pt}}
  \caption{Overview of the proposed pipeline. Darker cells are the key blocks retained by the selection pass.}
  \label{fig:overview}
\end{figure}

\s{We sweep the number of retained blocks from two to thirty two.}
\s{Accuracy saturates at sixteen blocks on every benchmark we tried.}
\s{Beyond that point the extra blocks carry almost no probability mass.}
\s{We therefore use sixteen retained blocks in all remaining experiments.}\footnote{The sweep for the two smallest benchmarks is reported in Appendix A.}

\s{We also vary the block size between sixteen and two hundred and fifty six positions.}
\s{Small blocks give a finer pattern but make the selection pass relatively more expensive.}
\s{Large blocks make the selection cheap and the pattern coarse.}
\s{A block size of sixty four is the best compromise on our hardware.}

\subsection{Ablations}

\s{Removing the diagonal block from the retained set costs four points of accuracy.}
\s{Removing the prefix block costs two points and mostly affects instruction following.}
\s{Replacing the learned selection with a random one costs eleven points.}
\s{This last number is the clearest evidence that the selection itself carries the improvement.}

\section{Discussion}

\s{The method is orthogonal to quantisation and to key value cache compression.}
\s{We expect the savings to multiply rather than overlap, but we have not measured the combination.}
\s{Our selection rule assumes that block summaries are informative about block contents.}
\s{This assumption fails when a block mixes several unrelated topics, which happens at document boundaries.}
\s{Inserting a separator token at those boundaries removes most of the degradation.}

\section{Conclusion}

\s{We presented a block sparse attention kernel that chooses its sparsity pattern from the data.}
\s{It matches dense attention on five long document benchmarks and runs three times faster at thirty two thousand tokens.}
\s{The remaining gap to the theoretical peak is dominated by the selection pass, which we expect to shrink with better summaries.}

\begin{thebibliography}{9}
\bibitem{doe2024sparse}
J. Doe and A. Smith. Sparse transformers for long context. In Proc. NeurIPS, 2024.
\bibitem{roe2025efficient}
R. Roe. Efficient attention in practice. Journal of ML, 12(3):45--67, 2025.
\bibitem{tanaka2023blocks}
M. Tanaka. Block structured kernels. In Proc. ICML, pages 991--1002, 2023.
\end{thebibliography}

% 付録を参考文献の **後ろ** に置く。ML 系の論文では普通の並びで、しかも
% 「参考文献より後ろは全部 reference」という切り方だとここが丸ごと消える。
\appendix

\section{Additional Sweeps}

\s{This appendix reports the sweeps that did not fit in the main text.}
\s{The smallest two benchmarks saturate at eight retained blocks rather than sixteen.}
\s{We kept sixteen everywhere so that a single configuration covers all five datasets.}
\s{The difference in accuracy between eight and sixteen blocks on those two benchmarks is under half a point.}

\section{Implementation Notes}

\s{The kernel is written for a single head and is launched once per head and per batch element.}
\s{Block summaries are computed in the same launch as the previous layer, so they add no extra pass over memory.}
\s{We keep the summaries in half precision and the running softmax statistics in single precision.}
\s{Mixing the two precisions was necessary to keep the numerical error below the noise between seeds.}

\end{document}
